\section{Описание}

В работе реализован поиск множества заранее неизвестных образцов в одном известном тексте. Так как
текст считывается первым и больше не меняется, по нему строится суффиксное дерево. После этого каждый
образец можно искать как путь от корня дерева: если путь существует, то все листья в соответствующем
поддереве задают позиции вхождений образца в текст.

\subsection{Построение дерева}

Суффиксное дерево строится онлайн-алгоритмом Укконена. В конце текста
добавляется терминальный символ \texttt{\$}, чтобы все суффиксы заканчивались в листьях. Рёбра дерева
не хранят строки явно: каждое ребро задаётся границами \texttt{beg} и \texttt{end} в общей строке.
Для листьев используется общий правый конец \texttt{global\_end}, поэтому при добавлении очередного
символа все листовые рёбра удлиняются неявно.

Каждая вершина содержит:
\begin{itemize}
\item \texttt{beg} и \texttt{end} -- границы подстроки на ребре, ведущем в вершину;
\item \texttt{idx} -- индекс начала суффикса для листа;
\item \texttt{children} -- переходы по первому символу ребра;
\item \texttt{suf\_link} -- суффиксную ссылку для алгоритма Укконена.
\end{itemize}

Переходы хранятся в \texttt{std::unordered\_map<char, Node*>}. Для данной задачи лексикографический
порядок переходов не нужен: позиции найденных вхождений после поиска сортируются отдельно.

\subsection{Расширение дерева}

Метод \texttt{extend} добавляет очередной символ текста. Рабочая точка алгоритма задаётся переменными
\texttt{active\_node}, \texttt{active\_edge} и \texttt{active\_length}. Счётчик \texttt{remaining}
показывает, сколько суффиксов текущей фазы ещё нужно обработать явно.

Внутри расширения используются стандартные случаи алгоритма Укконена:
\begin{itemize}
\item если из активной вершины нет нужного ребра, создаётся новый лист;
\item если символ на существующем ребре совпадает с добавляемым, фаза завершается увеличением
\texttt{active\_length};
\item если символ отличается, ребро разбивается, создаётся внутренняя вершина и новый лист.
\end{itemize}

После явного добавления суффикса рабочая точка обновляется: либо выполняется переход по суффиксной
ссылке, либо применяется специальное правило для корня.

\subsection{Поиск образца}

Метод \texttt{search} последовательно сопоставляет символы образца с рёбрами суффиксного дерева. Если
на каком-либо шаге нужного перехода нет или символ внутри ребра отличается, образец отсутствует в
тексте. Если образец полностью совпал с путём в дереве, выполняется обход поддерева и собираются
индексы всех листьев.

Так как при использовании \texttt{unordered\_map} порядок обхода детей не определён, перед выводом
позиции сортируются. В самой программе позиции переводятся из нумерации с нуля в нумерацию с единицы.

Асимптотика реализации:
\begin{itemize}
\item построение дерева~--- $O(n)$, где $n$~--- длина текста; алгоритм Укконена
      выполняет $O(n)$ операций в худшем случае, каждый переход по
      \texttt{unordered\_map} занимает $O(1)$ амортизированно в среднем;
\item поиск одного образца~--- $O(m + k \log k)$, где $m$~--- длина образца,
      $k$~--- число вхождений; слагаемое $O(m)$ соответствует спуску от корня,
      $O(k)$~--- обходу поддерева, $O(k \log k)$~--- сортировке позиций перед выводом;
\item память~--- $O(n)$: суффиксное дерево содержит не более $2n$ вершин
      (не более $n$ листьев и $n - 1$ внутренних вершин).
\end{itemize}

\section{Исходный код}

Ниже приведены основные фрагменты реализации.

\lstinputlisting[
language=C++,
firstline=9,
lastline=22,
caption={Структура вершины суффиксного дерева},
label={lst:node}
]{../main.cpp}

\lstinputlisting[
language=C++,
firstline=61,
lastline=77,
caption={Поля дерева и создание вершин},
label={lst:fields}
]{../main.cpp}

\lstinputlisting[
language=C++,
firstline=79,
lastline=140,
caption={Добавление символа алгоритмом Укконена},
label={lst:extend}
]{../main.cpp}

\lstinputlisting[
language=C++,
firstline=35,
lastline=57,
caption={Поиск образца в суффиксном дереве},
label={lst:search}
]{../main.cpp}

\lstinputlisting[
language=C++,
firstline=142,
lastline=148,
caption={Сбор позиций в листьях поддерева},
label={lst:collect}
]{../main.cpp}

\lstinputlisting[
language=C++,
firstline=152,
lastline=177,
caption={Чтение входа и вывод найденных позиций},
label={lst:main}
]{../main.cpp}
\pagebreak

\section{Консоль}
\begin{alltt}
$ g++ -std=c++17 -O2 -Wall -Wextra -pedantic main.cpp -o lab5.out
$ ./lab5.out
abcdabc
abcd
bcd
bc
1: 1
2: 2
3: 2, 6
\end{alltt}
\pagebreak
